\documentclass[11pt,a4paper]{article}
\usepackage[utf8]{inputenc}
\usepackage[T1]{fontenc}
\usepackage{amsmath,amssymb,amsthm}
\usepackage{booktabs,array,caption,microtype,xcolor}
\usepackage[margin=2.5cm]{geometry}
\linespread{1.06}
\setlength{\emergencystretch}{3em}
\allowdisplaybreaks[1]
\usepackage[colorlinks=true,linkcolor=blue!60!black,citecolor=blue!60!black,urlcolor=blue!60!black]{hyperref}
\newtheorem{theorem}{Theorem}
\newtheorem{proposition}[theorem]{Proposition}
\newtheorem{lemma}[theorem]{Lemma}
\newtheorem{conjecture}[theorem]{Conjecture}
\newtheorem{corollary}[theorem]{Corollary}
\theoremstyle{remark}\newtheorem{remark}[theorem]{Remark}
\newcommand{\Z}{\mathbb Z}\newcommand{\R}{\mathbb R}\newcommand{\Q}{\mathbb Q}\newcommand{\HH}{\mathbb H}
\newcommand{\e}{\mathrm e}\newcommand{\Per}{\mathrm{Per}}
\newcommand{\todo}[1]{\textcolor{red!70!black}{\small\sffamily[#1]}}
\newcommand{\status}[1]{\textcolor{blue!50!black}{\small\sffamily Status: #1}}

\title{Dilated pieces, Hecke correspondences and the second spectrum:\\ the level-one lines of the pieces $\mathcal P_u$}
\author{Jasper Reichardt\thanks{Independent researcher. Correspondence: jasperreichardt@gmail.com. This work was
carried out with heavy use of a large language model; see the disclosure at the end.}}
\date{Working draft of 15 September 2026 --- not for circulation. Every displayed constant was tested numerically before
it was written down; the proofs marked ``sketch'' are not complete.}

\begin{document}
\maketitle

\begin{abstract}
Part~IV of this series~\cite{ReichardtIV} showed that the pieces $\mathcal P_u(Y)$ of the off-diagonal of the pair
singular series of an irreducible quadratic $f$ oscillate at the spectral parameters of the even Maass cusp forms, and
that the piece with dilation $u$ lives at level $u^2$. It left open ``the size, in $u$, of the restricted geometric
factor''. This note computes that factor for the level-one lines and for the newforms of level $u$. The piece with
dilation $u$ is the level-one seed of the smooth-window theorem, at length $uY$, summed over the $\Gamma_0(u^2)$-invariant
family $\mathcal W_u$ of forms of discriminant $4u^2D$ with $u^2\mid a$ and $2u^2\mid b$. The coefficient of the line $t_j$ of
a level-one form $u_j$ in piece $u$ is therefore a projection onto the oldform space $\{u_j(z),u_j(uz),u_j(u^2z)\}$, and
for an odd prime $u$ it equals the coefficient in piece $1$ times an explicit rational function of $\lambda_j(u)$,
$\sqrt u$ and $\chi_D(u)$: the Gram matrix of the oldforms (Iwaniec--Luo--Sarnak) and two local period identities for the
Hecke neighbours of Heegner points (Gross--Kohnen--Zagier). The decay in $u$ is set by the splitting type of $u$ in
$\Q(\sqrt D)$ and not by the Hecke eigenvalue: the ratio is exactly $u^{-3/2}$ when $u$ is inert, $-u^{-3/2}(1-2\lambda_j(u)u^{-1/2}+O(u^{-1}))$
when $u$ splits, and $u^{-1/2}(1-\lambda_j(u)u^{-1/2}+O(u^{-1}))$ when $u$ is ramified; so the main terms of the level-one lines
are uniformly summable over $u$. The same formula with a newform $v$ of level $u$ gives its line in piece $u$, and the
period that drives it vanishes when $u$ is inert, because the Hecke and Fricke contributions cancel. Both predictions carry
no free parameter. Where the line is resolved in the data ($|r|\gtrsim0.1$ on grids to $Y=10^7$) they hold to a few percent
in amplitude and a few hundredths of a radian in phase --- $13$ of the $21$ pairs $(u,D)$ tested for the first level-one line,
with one unexplained outlier, and five of the seven discriminants for the first newform of level~$3$, the other two being the
inert cases, where the line is observed at the noise level as predicted. The error term uniform in $u$, the open core of the
conjecture of part~III, is not touched.
\end{abstract}

\section{Introduction}\label{sec:intro}

\paragraph{The object.} Notation as in part~IV~\cite{ReichardtIV}. For a discriminant $D<0$, a squarefree $u\ge1$ and
$Q_u(h)=u^2h^2-D$, the unrestricted piece with dilation $u$ and smooth window $w\in C_c^\infty((0,\infty))$ is
\[
S^w_u(Y)=\sum_{h\ge1}w\Bigl(\frac hY\Bigr)\bigl(\sigma_{-1}(Q_u(h))-\mathbb E_u\bigr),\qquad
\sigma_{-1}(n)=\sum_{d\mid n}\frac1d,\quad \mathbb E_u=\sum_d\frac{\rho_u(d)}{d^2},
\]
$\rho_u(d)$ the number of roots of $u^2x^2\equiv D\pmod d$. For $u=1$ Theorem~\textsc{smooth} of part~IV gives
\begin{equation}\label{eq:smooth1}
\begin{gathered}
\sqrt Y\,S^w_1(Y)=|D|^{-1/4}\sum_j\Per_D(u_j)\,\tilde L_j\bigl[A_j(Y)+\overline{A_j(Y)}\bigr]+\mathcal E^w(Y)+O(Y^{-2+\varepsilon}),\\
A_j(Y)=\Gamma(-it_j)W_c\bigl(\tfrac32+it_j\bigr)\Bigl(\frac{\pi\sqrt{|D|}}Y\Bigr)^{it_j},
\end{gathered}
\end{equation}
the sum over an orthonormal basis of even Maass cusp forms of $\mathrm{SL}_2(\Z)$, with $\Per_D(u_j)=\sum_Qu_j(z_Q)/|\Gamma_Q|$
over all classes of forms of discriminant $4D$, $\tilde L_j=\sum_ka_j(k)k^{-3/2}$, and $W_c$ the Mellin factor of the window.
Part~IV, Section~7, showed that piece $u$ lives at level $u^2$ and named the one missing ingredient: how the coefficients
depend on $u$. Its footnote on $u=3,5$ recorded that the level-one lines are weak there and that the lines of $\Gamma_0(9)$
take over.

\paragraph{Results.} Write $\lambda_j(u)$ for the Hecke eigenvalue ($T_uu_j=\lambda_j(u)u_j$, $a_j(1)=1$) and
$\chi=\chi_D(u)$ for the Kronecker symbol, so $\chi=0$ iff $u\mid D$.

\begin{theorem}[the spectral formula for piece $u$]\label{thm:pieceu}
Let $\mathcal W_u=\{[a,b,c]\ \text{of discriminant }4u^2D:\ u^2\mid a,\ 2u^2\mid b\}$ (all contents). Then $\mathcal W_u$ is
$\Gamma_0(u^2)$-invariant and, for every $\varepsilon>0$,
\[
\sqrt Y\,S^w_u(Y)=\sqrt u\,|D|^{-1/4}\sum_v\Per_{\mathcal W_u}(v)\,\tilde L_v\bigl[A_v(uY)+\overline{A_v(uY)}\bigr]+\mathcal E^w_u(Y)+O(Y^{-2+\varepsilon}),
\]
the sum over an orthonormal basis $\{v\}$ of the even Maass cusp forms of $\Gamma_0(u^2)$, with $\Per_{\mathcal W_u}(v)=
\sum_{Q\in\Gamma_0(u^2)\backslash\mathcal W_u}v(z_Q)/|\Gamma_0(u^2)_{z_Q}|$, $\tilde L_v$ formed with the Fourier coefficients
of $v$ at the cusp $\infty$, and $\mathcal E^w_u$ the continuous spectrum.
\end{theorem}
\status{proved (Section~\ref{sec:seed}): it is Theorem~\textsc{smooth} of part~IV at level $u^2$ applied to an explicit
identity between seeds; every constant checked numerically at $u=2,3,5$.}

\begin{theorem}[the level-one lines]\label{thm:level1}
Let $u$ be an odd prime and $u_j$ an even Hecke--Maass form of level one. The coefficient of the line $t_j$ in
$\sqrt Y S^w_u(Y)$ equals the coefficient in $\sqrt YS^w_1(Y)$ multiplied by $u^{-it_j}$ and by
\begin{equation}\label{eq:r}
r_j(u;D)=\sqrt u\,P^{\mathsf T}\bigl(u(u+1)\,G\bigr)^{-1}L,\qquad
G=\begin{pmatrix}1&g&h\\ g&1&g\\ h&g&1\end{pmatrix},
\end{equation}
\[
g=\frac{\lambda_j(u)\sqrt u}{u+1},\qquad h=\frac{\lambda_j(u)^2-1-1/u}{u+1},\qquad L=(1,u^{-1},u^{-2}),\qquad P=(p_0,p_1,p_0),
\]
\[
p_1=\begin{cases}2u&u\mid D,\\ u-\chi&u\nmid D,\end{cases}\qquad
p_0=\lambda_j(u)\sqrt u+\begin{cases}u-1&u\mid D,\\ -(1+\chi)&u\nmid D.\end{cases}
\]
Moreover
\[
\begin{aligned}
r_j(u;D)&=u^{-3/2}&&\text{(inert),}\\
r_j(u;D)&=-u^{-3/2}\bigl(1-2\lambda_j(u)u^{-1/2}+O(u^{-1})\bigr)&&\text{(split),}\\
r_j(u;D)&=u^{-1/2}\bigl(1-\lambda_j(u)u^{-1/2}+(\lambda_j(u)^2+1)u^{-1}+O(u^{-3/2})\bigr)&&\text{(ramified),}
\end{aligned}
\]
the implied constants absolute, and the inert identity exact for every $\lambda_j(u)$. The Hecke eigenvalue enters only the corrections; the decay is set by the splitting type.
\end{theorem}
\status{proved (Section~\ref{sec:lemmas}) for $u\nmid D$ and for $u\parallel D$ with $u$ not dividing the conductor of $D$;
the asymptotics are read off~\eqref{eq:r} symbolically (the inert identity was verified with a computer algebra system and
is an algebraic identity in $\lambda$ and $\sqrt u$). The first version of this draft stated the asymptotic as
$\lambda_j(u)/u$ for $u\nmid D$; that was wrong --- the $\lambda_j(u)\sqrt u$ terms of $p_0$ and $p_1L_2$ cancel exactly --- and was
caught by an outside reading that evaluated~\eqref{eq:r} at large $u$. The values of~\eqref{eq:r} agree with the direct
spectral computation to four digits and with the data to the noise (Table~\ref{tab:ratios}).}

\begin{theorem}[the newform lines]\label{thm:newform}
Let $u$ be an odd prime and $v$ an even Hecke--Maass newform of level $u$ with Atkin--Lehner eigenvalue $\varepsilon$
($a_v(u)=-\varepsilon u^{-1/2}$). The coefficient of its line in $\sqrt YS^w_u(Y)$ is
\[
2\sqrt u\,|D|^{-1/4}\Bigl[\sum_{k=1}^2\Per_{\mathcal W_u}(v_k)\tilde L(v_k)\Bigr]\,\Gamma(-it_v)W_c\bigl(\tfrac32+it_v\bigr)\,u^{-it_v},
\]
$\{v_1,v_2\}$ the orthonormalisation of $\{v(z),v(uz)\}$ on $\Gamma_0(u^2)\backslash\HH$, whose Gram matrix has off-diagonal
entry $a_v(u)\sqrt u/u=-\varepsilon/u$ relative to $\langle v,v\rangle$. If $u$ is inert in $\Q(\sqrt D)$ then
$\Per_{\mathcal W_u}(v)=\Per_{\mathcal W_u}(v(u\,\cdot))=0$ and the line is absent.
\end{theorem}
\status{proved: the formula is Theorem~\ref{thm:pieceu} restricted to the two-dimensional oldspace, the Gram entry is
Lemma~\ref{lem:gram}(iii), and the vanishing is Lemma~\ref{lem:newperiod}. Tested on seven discriminants (Table~\ref{tab:newform}):
the period vanishes identically for inert $u$ in the computation, and in the data the line is at the noise level there.}

\begin{corollary}[uniform summability of the level-one main terms]\label{cor:sum}
For a fixed level-one form $u_j$ and any bounded weights $w(u)$, $\sum_{u\le U}|w(u)|\,(H/u)^{1/2}\,|r_j(u;D)|\ll H^{1/2}$
uniformly in $U$: the sum over odd primes $u\nmid D$ converges absolutely like $\sum_uu^{-2}$, and the finitely many
$u\mid D$ contribute $\ll_D H^{1/2}$.
\end{corollary}

\begin{remark}[what this does and does not settle]
Part~III's conjecture needs the pieces with $u\le H^{2/3}$ to sum to $o(H\log H)$ with their \emph{errors}. The
corollary says the level-one \emph{main} terms do so with room to spare; the newform lines of level $u$ are new at every
$u$ (Theorem~\ref{thm:newform} gives their size for each $u$ but not a bound uniform in the newform), the cycloidal-group
forms of level $u^2$~\cite{Stromberg2012} are not covered at all, and the error term of Remark~\textsc{gaps} of part~IV is
untouched. What has changed is that the main-term side of the $u$-problem is Hecke theory on Heegner divisors, and it
decays.
\end{remark}

\paragraph{Method and checks.} Section~\ref{sec:seed} proves Theorem~\ref{thm:pieceu}. Section~\ref{sec:lemmas} states
the three lemmas behind Theorem~\ref{thm:level1} with what is proved of each. Section~\ref{sec:numerics} reports the
tests: the ratio~\eqref{eq:r} against a direct computation of the oldform projection with the Maass forms' values (no
closed form used), and both against the pieces computed to $Y=10^7$; then the newform of level~$3$; then the case $u=2$
and the composite $u=6$. Section~\ref{sec:open} lists what is open. Scripts and data are in the repository
(\texttt{research/explore/}, \texttt{hecke-oldform-predict.py}, \texttt{hecke-newform-u3.py}, \texttt{hecke-ratio-sharp.py}).

\section{The seed identity and the family $\mathcal W_u$}\label{sec:seed}

The pairs $(d,x)$ with $u^2x^2\equiv D\pmod d$ are the forms $Q=[d,\,2ux,\,(u^2x^2-D)/d]$ of discriminant $4D$ with
$2u\mid b$; write $\mathcal H^{(u)}$ for this sub-family of the set $\mathcal H$ of all forms of discriminant $4D$, and
$z_Q=(-ux+i\sqrt{|D|})/d$. Poisson summation in $h$, exactly as in the proof of Theorem~\textsc{smooth}, writes
$S^w_u(Y)$ as $\sum_{(d,x)}\frac{Y}{d^2}\sum_{k\ne0}\hat w(kY/d)\,\e(kx/d)$, and in the coordinates of $z=z_Q$ ($y=\sqrt{|D|}/d$,
$x/d=-\Re z/u$) the summand is
\[
\frac{Yy^2}{|D|}\sum_{k\ne0}\hat w\Bigl(\frac{kYy}{\sqrt{|D|}}\Bigr)\e\Bigl(-\frac{k\Re z}u\Bigr)=u\,\Psi^w_{uY}(z/u),
\]
$\Psi^w_{Y'}(\tau)=\frac{Y'(\Im\tau)^2}{|D|}\sum_{k\ne0}\hat w(kY'\Im\tau/\sqrt{|D|})\e(-k\Re\tau)$ the level-one seed of
part~IV at length $Y'$. The point $\tau=z_Q/u=(-2u^2x+i\sqrt{4u^2|D|})/(2u^2d)$ is the Heegner point of the form
$Q'=[u^2d,\,2u^2x,\,(u^2x^2-D)/d]$ of discriminant $4u^2D$, and $Q\mapsto Q'$ is a bijection of $\mathcal H^{(u)}$ onto
$\mathcal W_u$. For $\gamma\in\Gamma_0(u^2)$ the middle coefficient transforms as $b'=2a\alpha\beta+b(\alpha\delta+\beta\gamma)+2c\gamma\delta
\equiv b\,\alpha\delta\equiv b\pmod{2u^2}$, since $u^2\mid a$, $u^2\mid\gamma$ and $\alpha\delta\equiv1\pmod{u^2}$; the
condition $u^2\mid a$ is $\Gamma_0(u^2)$-stable by Lemma~\textsc{subfamily} of part~IV. So $\mathcal W_u$ is a finite union of
$\Gamma_0(u^2)$-orbits, the translation $\tau\mapsto\tau+1$ is $x\mapsto x-d$, and $S^w_u(Y)=u\sum_{\tau\in\mathcal W_u/\Z}\Psi^w_{uY}(\tau)$.
The proof of Theorem~\textsc{smooth} applies verbatim on $\Gamma_0(u^2)\backslash\HH$ (its cusp at $\infty$ has width one, so
the pairing with the Fourier expansions at $\infty$ is unchanged), with $Y$ replaced by $uY$; the factor $u$ and
$(uY)^{-1/2}$ combine to $\sqrt u\,Y^{-1/2}$, and $(\pi\sqrt{|D|}/(uY))^{it_v}$ carries the phase $-t_v\log u$. This proves
Theorem~\ref{thm:pieceu}. For $f_m(z)=u_j(u^mz)$ the coefficient at index $u^mn$ is $u^{m/2}a_j(n)$, so $\tilde L_{f_m}=u^{-m}\tilde L_j$,
which is the vector $L$ of Theorem~\ref{thm:level1}.

\section{The three lemmas}\label{sec:lemmas}

Throughout $u$ is an odd prime and $\Per_D(f)=\sum_{Q\in\mathcal H/\mathrm{SL}_2(\Z)}f(z_Q)/|\Gamma_Q|$.

\begin{lemma}[multiplicity]\label{lem:mult}
Let $u$ be an odd prime with $u\nmid D$, or $u\parallel D$ and $u$ not dividing the conductor of $D$. For every
$\mathrm{SL}_2(\Z)$-class $C$ of forms of discriminant $4D$,
$\sum_{Q'\in\Gamma_0(u^2)\backslash\mathcal W_u,\ uz_{Q'}\in C}|\Gamma_0(u^2)_{z_{Q'}}|^{-1}=m_u(D)\,|\Gamma_C|^{-1}$ with
$m_u(D)=2u$ if $u\mid D$ and $m_u(D)=u-\chi_D(u)$ if $u\nmid D$. Equivalently $\Per_{\mathcal W_u}(u_j(u\,\cdot))=m_u(D)\Per_D(u_j)$.
\end{lemma}
\begin{proof}
Write $(Q\circ\gamma)(v)=Q(\gamma v)$ (a right action), $C=Q_0\circ\Gamma$ with $\Gamma=\mathrm{SL}_2(\Z)$, $S=\{\gamma:Q_0\circ\gamma=Q_0\}$,
and let $B$ be the bilinear form of $Q_0$, $B(v,w)=av_1w_1+\tfrac b2(v_1w_2+v_2w_1)+cv_2w_2$, of discriminant $D$.

\emph{The transported group.} With $\alpha=\left(\begin{smallmatrix}1&0\\0&u\end{smallmatrix}\right)$, $z\mapsto z/u$, one has
$\alpha^{-1}\Gamma_0(u^2)\alpha=G_u:=\Gamma_0(u)\cap\Gamma^0(u)=\{u\mid b,\ u\mid c\}$, of index $u(u+1)$. The bijection
$Q\mapsto Q'=[u^2a,ub,c]$ of $\mathcal H^{(u)}$ onto $\mathcal W_u$ has $z_{Q'}=z_Q/u$ and intertwines the action of $G_u$ on
$\mathcal H^{(u)}$ with that of $\Gamma_0(u^2)$ on $\mathcal W_u$, orbits and stabilisers included; and $\mathcal H^{(u)}$ is
$G_u$-stable, since $b(Q\circ g)\equiv b\,\alpha\delta\equiv b\pmod u$ for $g\in G_u$. So we count $G_u$-orbits in
$C\cap\mathcal H^{(u)}$.

\emph{Cosets are ordered pairs of lines.} $\Gamma\to\mathrm{SL}_2(\mathbb F_u)$ is onto and $G_u$ is the preimage of the diagonal
torus $T$, so $\Gamma/G_u\cong\mathrm{SL}_2(\mathbb F_u)/T$: a matrix with columns $(v,w)$ modulo $(v,w)\mapsto(\lambda v,\lambda^{-1}w)$
is the ordered pair of distinct lines $(\ell_1,\ell_2)=(\langle v\rangle,\langle w\rangle)$ in $\mathbb P^1(\mathbb F_u)$,
$u(u+1)$ pairs in all.

\emph{The family condition is orthogonality.} The middle coefficient of $Q_0\circ\gamma$ is $2B(\gamma e_1,\gamma e_2)$; so
$Q_0\circ\gamma\in\mathcal H^{(u)}$ iff $\ell_1\perp_B\ell_2$, a condition on the right coset $\gamma G_u$ and unchanged under
$\gamma\mapsto s\gamma$, $s\in S$.

\emph{The mass formula.} $C\cong S\backslash\Gamma$, the $G_u$-orbits in $C$ are the double cosets $S\gamma G_u$, with stabiliser
$G_u\cap\gamma^{-1}S\gamma$, and each double coset consists of $|S|/|G_u\cap\gamma^{-1}S\gamma|$ right cosets $\gamma'G_u$. Hence for
every function $\varphi$ on $\Gamma$ that is left-$S$- and right-$G_u$-invariant,
\begin{equation}\label{eq:mass}
\sum_{\text{$G_u$-orbits }O\subset C}\frac{\varphi(O)}{|\mathrm{Stab}_{G_u}(O)|}=\frac1{|S|}\sum_{\gamma\in\Gamma/G_u}\varphi(\gamma).
\end{equation}
With $\varphi=\mathbf 1[\ell_1\perp\ell_2]$ this gives $m_u(D)=\#\{(\ell_1,\ell_2)\text{ distinct}:\ \ell_1\perp_B\ell_2\}$. If
$u\nmid D$, $B$ is nondegenerate modulo $u$; every $\ell_1$ has a unique $\ell_1^\perp$, admissible iff $\ell_1^\perp\ne\ell_1$, i.e.\ iff
$\ell_1$ is not isotropic; the isotropic lines are the roots of $Q_0$ modulo $u$, $1+\chi_D(u)$ of them; so $m=u+1-(1+\chi)=u-\chi$.
If $u\parallel D$ and $u$ does not divide the conductor, $Q_0\not\equiv0$ and $B$ has rank one modulo $u$ with radical line $r$
(the double root); $B(v,w)=0$ iff $\langle v\rangle=r$ or $\langle w\rangle=r$, giving $u+u=2u$ ordered pairs.
\end{proof}

\begin{lemma}[the period of the dilated family]\label{lem:period}
Under the hypotheses of Lemma~\ref{lem:mult},
$\Per_{\mathcal W_u}(u_j)=\Per_{\mathcal W_u}(u_j(u^2\,\cdot))=\bigl(\lambda_j(u)\sqrt u+c_u(D)\bigr)\Per_D(u_j)$ with
$c_u(D)=u-1$ if $u\mid D$ and $c_u(D)=-(1+\chi_D(u))$ if $u\nmid D$.
\end{lemma}
\begin{proof}
In the notation of the previous proof, $z_{Q_0\circ\gamma}=\gamma^{-1}z_0$, so $z_{Q_0\circ\gamma}/u=\alpha\gamma^{-1}z_0$ with
$\alpha\gamma^{-1}$ of determinant $u$. Since $\alpha M\alpha^{-1}=\left(\begin{smallmatrix}a&b/u\\ uc&d\end{smallmatrix}\right)$ is integral iff $u\mid b$, the coset
$\Gamma\alpha\gamma^{-1}$ depends on $\gamma$ only modulo right multiplication by $\Gamma^0(u)$, which fixes the second column
line $\ell_2=\langle\gamma e_2\rangle$ and moves the first; the map $\gamma\mapsto\Gamma\alpha\gamma^{-1}$ is a bijection of
$\Gamma/\Gamma^0(u)$ onto the $u+1$ cosets $\Gamma\backslash\Delta_u$ of the Hecke operator. (Check: $\gamma=\left(\begin{smallmatrix}1&-k\\0&1\end{smallmatrix}\right)$
has $\ell_1=\langle e_1\rangle$ fixed, $\ell_2=\langle(-k,1)\rangle$, and $\alpha\gamma^{-1}z=(z+k)/u$, $u$ distinct neighbours.)
Hence $u_j(\alpha\gamma^{-1}z_0)=:u_j(h_{\ell_2}(z_0))$ depends only on $\ell_2$, the points $h_\ell(z_0)$, $\ell\in\mathbb P^1(\mathbb F_u)$,
are the $u+1$ Hecke neighbours of $z_0$, and $\sum_\ell u_j(h_\ell(z_0))=\sqrt u\,\lambda_j(u)u_j(z_0)$. Applying~\eqref{eq:mass} with
$\varphi(\gamma)=u_j(\alpha\gamma^{-1}z_0)\mathbf 1[\ell_1\perp\ell_2]$ and summing over the classes,
\[
\Per_{\mathcal W_u}(u_j)=\sum_C\frac1{|S_C|}\sum_{\ell_2}N(\ell_2)\,u_j(h_{\ell_2}(z_C)),\qquad N(\ell_2)=\#\{\ell_1\ne\ell_2:\ \ell_1\perp\ell_2\}.
\]
If $u\nmid D$ then $N(\ell_2)=1$ for non-isotropic and $0$ for isotropic $\ell_2$, so the inner sum is
$\sqrt u\lambda_j(u)u_j(z_C)-\sum_{\ell\text{ isotropic}}u_j(h_\ell(z_C))$. The isotropic lines are the eigenlines of the order
$\mathcal O_{4D}$ acting on $L_C/uL_C$, $L_C=\Z+\Z z_C$, i.e.\ the $1+\chi$ sublattices $\mathfrak pL_C$ for the primes
$\mathfrak p\mid u$ of the order; their Heegner points have discriminant $4D$ again, with class $[\mathfrak p]\cdot C$
\cite[II]{GrossKohnenZagier1987}. Since $C\mapsto\mathfrak pC$ permutes the classes of each order and preserves $|S_C|$ (the
units of the order), $\sum_C|S_C|^{-1}\sum_{\text{iso}}u_j(h_\ell(z_C))=(1+\chi)\Per_D(u_j)$, and the claim follows.
If $u\parallel D$ with $u$ prime to the conductor, $N(r)=u$ and $N(\ell_2)=1$ for $\ell_2\ne r$, so the inner sum is
$\sqrt u\lambda_j(u)u_j(z_C)+(u-1)u_j(h_r(z_C))$, and $h_r(z_C)=z_{\mathfrak pC}$ for the ramified prime $\mathfrak p$, again a
permutation of the classes: $\Per_{\mathcal W_u}(u_j)=(\sqrt u\lambda_j(u)+u-1)\Per_D(u_j)$. Finally
$u^2z_{Q'}=u z_Q=\alpha'\gamma^{-1}z_0$ with $\alpha'=\left(\begin{smallmatrix}u&0\\0&1\end{smallmatrix}\right)$, whose Hecke coset is
determined by $\ell_1$; the counts $\#\{\ell_2\ne\ell_1:\ell_1\perp\ell_2\}$ equal the $N(\ell_2)$ by the symmetry of $B$, which gives
$\Per_{\mathcal W_u}(u_j(u^2\cdot))=\Per_{\mathcal W_u}(u_j)$. (The same equality follows from the Fricke involution $\tau\mapsto-1/(u^2\tau)$
of $\Gamma_0(u^2)$, which maps $[u^2a',2u^2b',c]$ to $[u^2c,-2u^2b',a']$ --- the pairing of the divisor $d$ with $Q_u(x)/d$ --- and
so permutes $\mathcal W_u$.)
\end{proof}

\begin{lemma}[the Gram matrix]\label{lem:gram}
On $\Gamma_0(u^2)\backslash\HH$, with $f_m(z)=u_j(u^mz)$: (i) $\langle f_0,f_0\rangle=u(u+1)\|u_j\|^2_{\mathrm{SL}_2(\Z)}$ and
$\langle f_1,f_0\rangle=\langle f_2,f_1\rangle=\frac{\lambda_j(u)\sqrt u}{u+1}\langle f_0,f_0\rangle$;
(ii) $\langle f_2,f_0\rangle=\frac{\lambda_j(u)^2-1-1/u}{u+1}\langle f_0,f_0\rangle$; (iii) for a newform $v$ of level $u$ with Fricke
eigenvalue $\varepsilon$, $\langle v(u\,\cdot),v\rangle=-\frac\varepsilon u\langle v,v\rangle$.
\end{lemma}
\begin{proof}
Write $\mathrm{Tr}_H^G F=\sum_{\gamma\in H\backslash G}F\circ\gamma$ for $H\subset G$ of finite index, so that
$\langle F,\Phi\rangle_H=\langle\mathrm{Tr}_H^GF,\Phi\rangle_G$ for $G$-invariant $\Phi$, and $\langle\Phi,\Psi\rangle_H=[G:H]\langle\Phi,\Psi\rangle_G$
for $G$-invariant $\Phi,\Psi$. (i) The cosets $\Gamma_0(u)\backslash\Gamma$ are $\left(\begin{smallmatrix}1&0\\ b&1\end{smallmatrix}\right)$,
$b\bmod u$, and $S$, and the points $u\gamma z$ over them are the $u+1$ Hecke neighbours of $z$; so
$\mathrm{Tr}_{\Gamma_0(u)}^\Gamma(u_j(u\,\cdot))=\sqrt u\lambda_j(u)u_j$, $\langle u_j(uz),u_j\rangle_{\Gamma_0(u)}=\sqrt u\lambda_j(u)\|u_j\|^2_\Gamma=
\frac{\sqrt u\lambda_j(u)}{u+1}\|u_j\|^2_{\Gamma_0(u)}$ \cite[Lemma~2.4]{IwaniecLuoSarnak2000}, and the ratio persists on $\Gamma_0(u^2)$
since both functions are $\Gamma_0(u)$-invariant; $\langle f_2,f_1\rangle$ reduces to the same by $z\mapsto z/u$, which carries
$\Gamma_0(u^2)\backslash\HH$ to $(\Gamma_0(u)\cap\Gamma^0(u))\backslash\HH$ of the same index. (ii) Trace from $\Gamma_0(u^2)$ to $\Gamma_0(u)$
over $\gamma_b=\left(\begin{smallmatrix}1&0\\ ub&1\end{smallmatrix}\right)$: $f_2(\gamma_bz)=u_j\bigl(\left(\begin{smallmatrix}u&0\\ b&1\end{smallmatrix}\right)uz\bigr)$,
and $\left(\begin{smallmatrix}u&0\\ b&1\end{smallmatrix}\right)\in\Gamma\left(\begin{smallmatrix}1&k\\0&u\end{smallmatrix}\right)$ with $k\equiv b^{-1}$ for
$b\ne0$, so $f_2(\gamma_bz)=u_j((uz+k)/u)$ for $b\ne0$ and $u_j(u^2z)$ for $b=0$; by the Hecke relation at the point $uz$,
$\mathrm{Tr}\,f_2=\sum_{k\bmod u}u_j((uz+k)/u)-u_j(z)+u_j(u^2z)=\sqrt u\lambda_j(u)u_j(uz)-u_j(z)$. Hence
$\langle f_2,f_0\rangle_{\Gamma_0(u^2)}=\sqrt u\lambda_j(u)\langle f_1,f_0\rangle_{\Gamma_0(u)}-\|u_j\|^2_{\Gamma_0(u)}
=\bigl(\tfrac{\lambda_j(u)^2u}{u+1}-1\bigr)\|u_j\|^2_{\Gamma_0(u)}=\tfrac{\lambda_j(u)^2-1-1/u}{u+1}\langle f_0,f_0\rangle$.
(iii) The same trace gives $\mathrm{Tr}\,v(u\,\cdot)=\sum_b v\bigl(\left(\begin{smallmatrix}1&0\\ b&1\end{smallmatrix}\right)uz\bigr)$, and
$\sum_bv\bigl(\left(\begin{smallmatrix}1&0\\ b&1\end{smallmatrix}\right)w\bigr)=-v(-1/w)$ because a newform has vanishing trace to
level one; with $w=uz$, $\mathrm{Tr}\,v(u\,\cdot)=-v(-1/(uz))=-\varepsilon v(z)$, so $\langle v(uz),v\rangle_{\Gamma_0(u^2)}=-\varepsilon\|v\|^2_{\Gamma_0(u)}=-\frac\varepsilon u\|v\|^2_{\Gamma_0(u^2)}$.
\end{proof}

\begin{lemma}[the newform periods]\label{lem:newperiod}
Let $v$ be a newform of level $u$ with Fricke eigenvalue $\varepsilon$, under the hypotheses of Lemma~\ref{lem:mult}. Then
$\Per_{\mathcal W_u}(v(u\,\cdot))=\varepsilon\Per_{\mathcal W_u}(v)$, and
\[
\Per_{\mathcal W_u}(v)=\begin{cases}0&\text{$u$ inert,}\\ -\sum_C|S_C|^{-1}\sum_{\ell_2\text{ isotropic}}v(h_{\ell_2}(z_C))&\text{$u$ split,}\\
(u-1)\sum_C|S_C|^{-1}v(h_r(z_C))&\text{$u\parallel D$,}\end{cases}
\]
where the $h_\ell(z_C)$ are taken as points of $X_0(u)$: the sums on the right are periods of $v$ over the Heegner points of
level $u$ and discriminant $4D$ above each class, in the sense of Gross--Kohnen--Zagier.
\end{lemma}
\begin{proof}
For fixed $\ell_2$ the $u$ choices of $\ell_1$ are $\gamma\mapsto\gamma\left(\begin{smallmatrix}1&0\\ c&1\end{smallmatrix}\right)$, and
$\alpha\left(\begin{smallmatrix}1&0\\-c&1\end{smallmatrix}\right)\alpha^{-1}=\left(\begin{smallmatrix}1&0\\-uc&1\end{smallmatrix}\right)\in\Gamma_0(u)$;
so $v(\alpha\gamma^{-1}z_0)=:V(\ell_2)$ depends only on $\ell_2$ also for the level-$u$ form. As points of $X_0(u)$ the $u+1$
values are $(z_0+k)/u$ ($\ell_2=\langle(-k,1)\rangle$) and $-1/(uz_0)$ ($\ell_2=\langle e_1\rangle$), so
$\sum_{\ell_2}V(\ell_2)=\sqrt u\,a_v(u)v(z_0)+\varepsilon v(z_0)=(-\varepsilon+\varepsilon)v(z_0)=0$: for a newform the Hecke term and the Fricke
term cancel. By~\eqref{eq:mass}, $\Per_{\mathcal W_u}(v)=\sum_C|S_C|^{-1}\sum_{\ell_2}N(\ell_2)V_C(\ell_2)$ with $N$ as in Lemma~\ref{lem:period};
for inert $u$, $N\equiv1$ and the sum vanishes; for split $u$, $N=1-\mathbf 1[\text{isotropic}]$; for $u\parallel D$, $N(r)=u$ and
$N=1$ otherwise. The relation between $v$ and $v(u\,\cdot)$ is the Fricke involution of $\Gamma_0(u^2)$ acting on $\mathcal W_u$ as the
divisor pairing, since $v(-1/(u^2\tau))=\varepsilon v(u\tau)$. Numerically the split and ramified formulas reproduce the direct orbit
sums to all computed digits at $D=-8,-11,-20$ and $D=-3,-15$ ($u=3$).
\end{proof}

\begin{proof}[Proof of Theorem~\ref{thm:level1} from the lemmas]
Theorem~\ref{thm:pieceu} restricted to the span of $f_0,f_1,f_2$ gives the coefficient $\sqrt u\sum_k\Per_{\mathcal W_u}(v_k)\tilde L(v_k)$
over an orthonormal basis $v_k$ of that span, which is the basis-free quantity $\sqrt u\,\mathbf P^{\mathsf T}\mathbf G^{-1}\mathbf L$
with $\mathbf P_m=\Per_{\mathcal W_u}(f_m)$, $\mathbf G_{mn}=\langle f_m,f_n\rangle$, $\mathbf L_m=\tilde L_{f_m}$. Dividing by
the $u=1$ coefficient $\Per_D(u_j)\tilde L_j/\|u_j\|^2$ and inserting the lemmas gives~\eqref{eq:r}; the phase $u^{-it_j}$
is from Theorem~\ref{thm:pieceu}. The asymptotics are not read off termwise --- $p_1\asymp u$ is large while $L_2=1/u$ and
$g\asymp\lambda u^{-1/2}$ are small, and the $\lambda_j(u)\sqrt u$ of $p_0$ cancels against $p_1(1/u-g)$ --- but from the closed
form~\eqref{eq:r} as a rational function of $\lambda$ and $\sqrt u$: with a computer algebra system, $r=u^{-3/2}$ identically
when $\chi=-1$, and the expansions in $u^{-1/2}$ stated in the theorem when $\chi=+1$ and $\chi=0$.
\end{proof}

\section{Numerical confirmation}\label{sec:numerics}

All data are the pieces of part~IV computed to $Y=10^7$ (\texttt{piece-divset.ts}, unrestricted object), sharp Riesz means
on $4096$ log-spaced points and smooth windows on $401$; amplitudes and phases by joint least squares at the level-one
lines with the dominant level-$u^2$ lines as nuisance frequencies (Str\"omberg's $\Gamma_0(9)$ eigenvalues at $u=3$
\cite{Stromberg2012}, the six strongest data-driven peaks at $u=5$). The predictions of Theorem~\ref{thm:level1} were
computed twice: from~\eqref{eq:r}, and directly by evaluating the Maass form at the orbit representatives of $\mathcal W_u$,
integrating the Gram matrix over the cosets of $\Gamma_0(u^2)$, and orthonormalising; the two agree to four digits, and
the direct route also checks $\langle f_0,f_0\rangle/\|u_j\|^2=u(u+1)$ and the Iwaniec--Luo--Sarnak entry.

\begin{table}[htbp]\centering\small
\caption{The first even line $t_1=13.7798$: signed ratio of its amplitude in piece $u$ to piece $1$ (sign $-$ when the
phase is shifted by $\pi$ beyond $-t_1\log u$). ``pred'' from~\eqref{eq:r} for odd $u$ and from the direct computation
for $u=2$; ``obs'' from the sharp grids. Predicted $|r|<0.1$ is below the noise of these grids.}\label{tab:ratios}
\begin{tabular}{@{}lrrrl@{}}
\toprule
$u$, type of $u$ in $\Q(\sqrt D)$ & pred & obs & $D$\\
\midrule
$2$, $D\equiv0\ (4)$, $D/4\equiv2,3\ (4)$ & $+0.850$ & $+0.852,\ +0.865,\ +0.899$ & $-4,-8,-20$\\
$2$, $D\equiv0\ (4)$, $D/4\equiv1\ (4)$ & $+0.945$ & (piece $u=2$ not yet fitted) & $-12$\\
$2$, split ($D\equiv1\ (8)$) & $-0.341$ & $-0.353,\ -0.379,\ -0.428$ & $-7,-15,-23$\\
$2$, inert ($D\equiv5\ (8)$) & $-0.068$ & $-0.065,\ -0.061,\ -0.071$ & $-3,-11,-19$\\
$3$, ramified & $+0.695$ & $+0.655,\ +0.726$ & $-3,-15$\\
$3$, split & $-0.311$ & $-0.285,\ -0.320,\ -0.332$ & $-8,-11,-20$\\
$3$, inert & $+0.192$ & $+0.196,\ (+0.341)$ & $-4,-7$\\
$5$, ramified & $+0.421$ & $+0.431$ & $-15$\\
$5$, split & $-0.076$ & (noise) & $-4,-11$\\
$5$, inert & $+0.089$ & (noise) & $-3,-8$\\
\bottomrule
\end{tabular}
\end{table}

\begin{table}[htbp]\centering\small
\caption{The first even newform of level $3$ ($t=5.0987$, Fricke $+1$) in the smooth piece $u=3$: predicted amplitude and
phase from Theorem~\ref{thm:newform} with the LMFDB coefficients, against the fit (all $\Gamma_0(9)$ even lines and four
level-one lines in the fit).}\label{tab:newform}
\begin{tabular}{@{}rlrrrr@{}}
\toprule
$D$ & $3$ in $\Q(\sqrt D)$ & $C_{\mathrm{pred}}$ & $C_{\mathrm{obs}}$ & ratio & $\Delta\varphi$\\
\midrule
$-3$ & ramified & $0.3200$ & $0.3262$ & $1.020$ & $-0.008$\\
$-15$ & ramified & $0.2016$ & $0.1946$ & $0.965$ & $+0.024$\\
$-8$ & split & $0.1404$ & $0.1426$ & $1.016$ & $+0.072$\\
$-11$ & split & $0.1371$ & $0.1370$ & $1.000$ & $+0.050$\\
$-20$ & split & $0.1945$ & $0.1931$ & $0.993$ & $+0.032$\\
$-4$ & inert & $0$ & $0.019$ & --- & ---\\
$-7$ & inert & $0$ & $0.006$ & --- & ---\\
\bottomrule
\end{tabular}
\end{table}

\paragraph{The case $u=2$.} The local constants are not integers: $\Per_{\mathcal W_2}(u_1(2\cdot))/\Per_D(u_1)=4.912$ for
$D/4\equiv2,3\pmod4$ and $5.278$ for $D/4\equiv1\pmod4$, $0.321$ for $D\equiv1\pmod8$, $1.373$ for $D\equiv5\pmod8$; so the
multiplicity of Lemma~\ref{lem:mult} is not class-independent at $u=2$ and the closed form~\eqref{eq:r} does not apply.
The direct computation does, and matches the data (Table~\ref{tab:ratios}). \todo{2-adic local analysis.}

\paragraph{Composite $u$.} For squarefree $u=u_1u_2$ the family condition factors over the primes and the oldform space is
$\{u_j(dz):d\mid u^2\}$, of dimension $9$ for $u=6$. The direct computation at level $36$ (Gram matrix over the $72$ cosets,
$\|f_0\|^2/\|u_j\|^2=72.000$) gives $r_1(6;-12)=+0.6573$ and $r_1(6;-8)=-0.2642$, against the products
$r_1(2;-12)\,r_1(3;-12)=0.945\times0.695=0.657$ and $r_1(2;-8)\,r_1(3;-8)=0.850\times(-0.311)=-0.264$; and the $9\times9$
Gram matrix is the tensor product of the $3\times3$ matrices at levels $4$ and $9$ entry by entry ($0.0781=0.7303\times0.1069$,
$-0.2323=0.7303\times(-0.3181)$, $-0.0955=0.3001\times(-0.3181)$, \dots). So the projection factors over the primes dividing
$u$ and the multiplicativity $r_j(u;D)=\prod_{p\mid u}r_j(p;D)$ holds on the theory side; its proof is the observation that the
period vector and the Gram matrix are tensor products, which follows from Lemmas~\ref{lem:mult}--\ref{lem:gram} applied
prime by prime. \todo{Write out the tensor-product argument.} The data at $u=6$ are noisy (the level-$36$ spectrum is dense):
the five signs observed agree with the products and the magnitudes scatter by up to $0.2$ ($D=-8,-20,-15,-4,-12$: observed
$-0.449,-0.320,-0.295,+0.111,+0.549$ against $-0.264,-0.264,-0.237,+0.163,+0.657$).

\paragraph{The cycloidal lines.} At $u=3$ the pieces are dominated by the lines $3.536$, $5.504$, $6.647$, $7.432$, $8.698$,
which Str\"omberg attributes to forms of the cycloidal group $\Gamma^3$~\cite{Stromberg2012}; their coefficients are not
tabulated and their amplitudes in piece $3$ cannot be predicted here.

\section{What is open}\label{sec:open}
\begin{enumerate}
\item The case $u$ dividing the conductor of $D$ (including $u=2$ for every $D$, since $4D$ has conductor divisible by $2$ when
$D$ is odd and $2$ ramified otherwise); the tensor-product argument for composite $u$; a proof of the inert identity
$r=u^{-3/2}$ that explains it rather than verifies it.
\item The $2$-adic case; the tensor-product argument for composite $u$.
\item The lines of the cycloidal-group forms, and any bound for the newform lines uniform in the newform: the level-aspect
equidistribution of \cite{LiuMasriYoung2013,HumphriesNordentoft2025}, route 1b of the knowledge base.
\item The error term uniform in $u$ (part~IV, Remark~\textsc{gaps}), which is the conjecture of part~III itself.
\end{enumerate}

\section*{AI-assisted research: disclosure}
This work was carried out with heavy use of large language models, the Anthropic models Claude Fable~5.1 and Claude
Opus~5, used interactively by the author in September 2026 as research assistants. The author posed the questions, set
the direction and the scope, and made every decision about what to pursue, claim and submit. The model derived the
identities, wrote and ran the software behind every number, searched the literature and drafted the text, in dialogue
with the author, who verified the code and the numerical claims against the repository record. The repository
(\url{https://github.com/gewure/ulamnd}) contains the scripts, the data, and the knowledge base with the record of errors.

\bibliographystyle{plain}
\bibliography{refs}
\end{document}
